\section{Hyperbolic ascending tori}\label{sec:02-ascending}

For an injective endomorphism $\phi:F\to F$, write
\[
 J_\phi=\langle F,t\mid t^{-1}gt=\phi(g)\ (g\in F)\rangle.
\]
Its height homomorphism sends $t$ to $1$ and $F$ to $0$.
Changing the sign of height, or replacing the stable letter by its inverse,
will not affect any assertion below. A \emph{free splitting tree} means a
minimal simplicial tree with trivial edge stabilizers and finite quotient;
its vertex stabilizers are free factors. A path map of trees is
\emph{alignment preserving} if it maps each segment monotonically onto a
segment, allowing constant subsegments.

\begin{theorem}\label{asc:vcs}
Let $F$ be a free group of finite rank and let $\phi:F\to F$ be injective.
If $J_\phi$ is word-hyperbolic, then $J_\phi$ is virtually compact special.
\end{theorem}

We prove the theorem by induction on $\rk F$. The rank-zero group is
infinite cyclic. Rank one gives $\mathrm{BS}(1,n)$, with $n\ne0$, and is
incompatible with hyperbolicity. When $\phi$ is an automorphism, the
conclusion follows from the cubulation theorem for hyperbolic
free-by-cyclic groups \cite{HagenWise2016} and Agol's theorem
\cite{Agol2013}. For a nonsurjective injection, we either reduce rank or
construct a noncontracting map of a free splitting tree. In the expanding
case, its associated real-tree action gives a finite family of point
stabilizers with nontrivial height-zero part. They are virtually
ascending tori of smaller rank and will serve as peripherals.
Filling their height-zero parts makes the base virtually free; a finite
graph cover then separates the expanding map into irreducible train-track
blocks. We remove any failure of fundamental-group injectivity in those
blocks before applying Hagen--Wise. This supplies the good fillings for
Theorem~\ref{fill:criterion}. The nonexpanding case is handled directly
by a quasiconvex hierarchy.

The ascending result is also an input to the filling construction in
Section~\ref{sec:06-closure}. After proving it, we recover its stronger
virtual finite-rank free-by-$\Z$ consequence in
Corollary~\ref{cor:ascending-finite-rank}, including the induction branches
that do not directly use the filling criterion.

\subsection{Finiteness hypotheses and elementary changes of model}

\begin{lemma}\label{asc:field-hypotheses}
An ascending torus of an injective finite-rank free-group endomorphism is
torsion-free, locally indicable, coherent, residually finite, and has a
finite aspherical complex of dimension at most two. It has the
Hughes-free subgroup-field setup used in Theorem~\ref{fill:criterion}, and
its finite cellular chain complex becomes acyclic over its rational
Linnell field.
\end{lemma}

\begin{proof}
The height kernel is an increasing union of conjugates of $F$.
Consequently a finitely generated subgroup contained in that kernel lies
in a free group and, if nontrivial, maps onto $\Z$. A subgroup not contained
in the kernel has a nonzero height homomorphism. This proves local
indicability and also torsion-freeness. Coherence is the theorem of
Feighn--Handel, which includes injective endomorphisms
\cite[Theorem~1.1]{FeighnHandel1999}. Residual finiteness is due to Borisov--Sapir
\cite{BorisovSapir2005}.

Represent $\phi$ by a based map of a finite rose. Its mapping torus is a
finite aspherical two-complex: the associated graph of spaces has
contractible vertex and edge universal covers and injective attaching
homomorphisms. The usual tree-of-spaces construction therefore gives a
contractible universal cover.

The height kernel belongs to Linnell's class generated by free groups
under directed unions; its extension by $\Z$ belongs to the same class
under elementary-amenable extensions. Linnell's theorem, together with
the locally indicable Hughes-free field setup, gives the rational
subgroup fields, their compatible inclusions and coset independence
\cite{Linnell1993,Hughes1970}; the locally indicable affiliated realization
and compatible rational subfields are given by
\cite[Corollaries~6.2--6.3]{JaikinLopezAlvarez2020} and
\cite[Section~2.3 and Proposition~2.2]{JaikinUniversality2021}.
In particular the field embeds in the
height skew-series division ring with coefficients in the kernel field,
with either order on height.

For completeness, take a basis $x_1,\ldots,x_r$ of $F$ and the relators
$t^{-1}x_it\phi(x_i)^{-1}$. The horizontal Fox matrix is
\[
 M=t^{-1}I-\left(\frac{\partial\phi(x_i)}{\partial x_j}\right)_{i,j}.
\]
After removing the invertible initial factor $t^{-1}$, its initial
height term is $I$. Geometric-series inversion therefore makes $M$
invertible in a height skew-series division ring. A square matrix over a
division subring that is invertible over an extension division ring is
already invertible over the subring, so $M$ is invertible over the
Linnell field. Thus the second boundary is injective after base change.
The zeroth homology vanishes because $t-1$ is invertible in that field.
The cellular Euler characteristic is $1-(r+1)+r=0$, so the first homology
vanishes as well.
\end{proof}

We will use the following elementary observation twice. It applies even
when a graph map is not injective on fundamental groups.

\begin{lemma}\label{asc:shift-equivalence}
Suppose based maps $f:X\to X$ and $g:Y\to Y$ of connected finite graphs
admit based maps $q:X\to Y$ and $i:Y\to X$ such that
\[
 qf=gq,\qquad ig=fi,\qquad iq=f^m,\qquad qi=g^m
\]
for some $m\ge1$. Their mapping-torus groups are isomorphic, by an
isomorphism preserving height.
\end{lemma}

\begin{proof}
The commuting maps induce homomorphisms of the mapping-torus
presentations, sending stable letter to stable letter. Their composites
are induced by $f^m$ and $g^m$. In a mapping-torus presentation the map
induced by the defining endomorphism is conjugation by the stable letter;
hence both composites are inner automorphisms. The two homomorphisms are
therefore isomorphisms. The same argument works up to based homotopy,
with the corresponding change of stable-letter marking.
\end{proof}

\subsection{An expanding relative tree model}

Expanding relative immersions for injective, nonsurjective free-group
endomorphisms
were constructed by Mutanguha
\cite[Proposition~4.1 and Theorem~4.5]{Mutanguha2021Dynamics}.
The following form is tailored to the induction on rank.

\begin{lemma}\label{asc:relative-model}
Let $\phi:F\to F$ be injective and nonsurjective, with $\rk F\ge2$.
After replacing $F$ by an iterated image, one of the following holds.
\begin{enumerate}[label=\textup{(\roman*)}]
\item A finite-index height subgroup of $J_\phi$ is an ascending torus of
an injection of a free group of rank strictly smaller than $\rk F$.
\item There is a nontrivial free splitting tree $T$ and a
$\phi$-equivariant embedding $f:T\to T$ that is a homothety of scale
$\lambda\ge1$. The edge-transition matrix is irreducible, every vertex
stabilizer is a proper free factor, and the quotient graph is finite.
\end{enumerate}
\end{lemma}

\begin{proof}
We compare successive iterated-image core graphs using the lengths of
their arcs between branch vertices. Collapsing the arcs whose lengths
stay bounded gives two trees and two maps: substitution identifies the
collapsed trees, while their original inclusion gives a possibly
collapsing map back. Their composition yields the required self-map.
Its transition diagram will either lose an edge orbit, allowing rank
reduction, or supply a recurrent irreducible part.

Fix the unit-edge Cayley tree $S$ of $F$. Put $F_n=\phi^n(F)$ and let
$S_n\subset S$ be its convex minimal subtree. Each $S_n$ is leafless.
Suppress the vertices of valence two and call the resulting edges
\emph{natural arcs}. The finite core graph $F_n\backslash S_n$ has rank
$r=\rk F$, and thus at most $3r-3$ natural edges.

Substitution by the fixed injection $\phi$ is a quasi-isometric embedding
of $S$: its image group is finitely generated and undistorted in $F$.
Restricting to $S_n$ and straightening to the hull $S_{n+1}$ gives
$\phi$-equivariant quasi-isometries
\[
 S_n\longrightarrow S_{n+1}
\]
with constants independent of $n$. Indeed the boundary of the target
hull is the image of the source boundary, and straightened axes remain
within the uniform quasigeodesic-stability bound. Branch points are
centers of tripods with three unbounded legs. These quasi-isometries
therefore give uniformly controlled branch-point correspondences in
both directions; the two composites move branch points a uniformly
bounded distance.

Choose a subsequence of pairs
\[
 U_k=S_{n_k},\qquad V_k=S_{n_k+1}
\]
on which the two finite reduced core-graph types are fixed and each of
their natural-edge lengths is either bounded or tends to infinity.
There is one bound $B$ for all bounded coordinates, and the minimum
$\ell_k$ of all unbounded coordinates tends to infinity. Collapse in
$U_k$ and $V_k$ the components of their respective bounded-arc forests.
These forests need not have bounded-diameter components.

For sufficiently large $k$, the branch-point correspondence induces a
$\phi$-equivariant simplicial isomorphism between the two collapsed
trees. Here are the details needed when a collapsed component is
unbounded. Adjacent branch points joined by a bounded arc have images
at a uniformly bounded distance. Their images consequently lie in the
same bounded-arc component, since that distance is less than $\ell_k$.
Apply this to chains of bounded arcs and then apply the coarse inverse;
the bounded errors of the two composites show that the resulting maps
on components are inverse bijections. To verify adjacency, let $[v,w]$
be a long natural source arc and choose the corresponding target branch
points $v',w'$. For every branch point $z$ of $[v',w']$, quasigeodesic
stability puts its corresponding source branch point within a uniform
distance of $[v,w]$. The projection of a branch point to the interior of
a branch-free natural arc is impossible: an off-arc branch point would
create a branch point where its joining arc meets the interior. Its
projection is therefore $v$ or $w$. It follows that $z$ is uniformly
close to $v'$ or $w'$. If $[v',w']$ crossed two long target arcs, a branch
point between them would be at distance at least $\ell_k$ from both
endpoints, a contradiction. It crosses at least one long arc because
the endpoint short components are distinct. Thus it crosses exactly
one. The reverse correspondence gives the same conclusion in the
opposite direction, proving adjacency.

The collapsed tree is nontrivial for some such subsequence. Otherwise
all natural arcs have uniformly bounded lengths along an infinite
subsequence. There would then be only finitely many unbased labelled
core graphs, so two nested groups $F_n,F_{n+j}$ would be conjugate. A
finitely generated subgroup of a free group cannot be conjugate to a
proper subgroup of itself. To see this last fact directly, if
$gAg^{-1}\le A$, the nested translates $g^{-i}S_A$ all contain $S_A$.
Fix a point on an axis in $S_A$. Since $A$ acts cocompactly on $S_A$,
that point is a uniformly bounded distance from every coset $g^{-i}A$.
Only finitely many cosets can meet a fixed finite ball, so $g^j\in A$
for some $j>0$. The nested conjugacy chain is now periodic and all its
inclusions are equalities. Hence $F_n=F_{n+j}$, which by injectivity
implies $F=\phi^j(F)$ and contradicts nonsurjectivity.

Inclusion $V_k\subset U_k$ maps bounded-arc components to bounded-arc
components. In fact a bounded natural arc of $V_k$ is subdivided by
branch points of $U_k$ into arcs of length at most $B<\ell_k$.
Fix one sufficiently large $k$, and write $n=n_k$ and
$T_U,T_V$ for the collapsed trees. Their actions are the actual actions
of $F_n$ and $F_{n+1}$, respectively. The maps just constructed are
\[
 T_U\xrightarrow{\ \Phi_k\ }T_V
     \xrightarrow{\ I_k\ }T_U.
\]
Here $\Phi_k$ is a simplicial isomorphism equivariant for
$\phi:F_n\xrightarrow{\sim}F_{n+1}$, and $I_k$ is alignment preserving
and equivariant for the inclusion $F_{n+1}\leq F_n$. Hence
$I_k\Phi_k$ is a $\phi$-equivariant alignment-preserving self-map of the
single $F_n$-tree $T_U$. The subsequence was used to classify natural-arc
lengths; no identification of the markings at different $k$ is needed.
Replacing the base by $F_n$
does not change the ascending torus: the base $F$ is conjugate into
$F_n$ by a power of the stable letter, and ascending normal form
gives the converse inclusion. Rename this base $F$.

If an iterate of the whole tree misses an edge orbit, its connected
image lies in one component of the complement of that orbit. The image
group $\phi^m(F)$ stabilizes that component. Its stabilizer $C$ is a
proper free factor, as follows by collapsing every complementary
component in the original free splitting. Consequently
\[
 \phi^m(F)\le C,\qquad \phi^m(C)\le C.
\]
The ascending torus of $\phi^m|_C$ embeds by normal form in the height-$m$
subgroup, and contains its whole base after conjugating by the stable
letter. It is therefore that subgroup. This is alternative (i).

Otherwise choose a source strongly connected component in the finite
edge-transition diagram. If it were nonrecurrent, it would consist of
one vertex with no self-loop and no incoming arrow; that edge orbit
would already be missing from the image of the whole tree. The chosen
component is therefore recurrent. No incoming arrows means that each
edge outside it maps entirely through edges outside it. Their union is
an invariant forest in the tree, so collapsing it gives a well-defined
self-map. Each remaining edge has a transition inside the recurrent
component and is not collapsed. Alignment preservation also survives:
lift a segment through
the collapsed vertex components, apply alignment preservation before
collapse, and delete the collapsed subsegments. No reduced path can
acquire backtracking in this operation. The transition matrix is now
irreducible. Its Perron eigenvector assigns positive lengths to the
finitely many edge types, and linear parametrization makes the map
multiply all path lengths by the Perron eigenvalue $\lambda\ge1$.
Alignment preservation and noncollapse make it an embedding.
All vertex stabilizers are proper free factors because the resulting
free splitting is nontrivial. This proves alternative (ii).
\end{proof}

\subsection{Point stabilizers and their quasiconvexity}

Assume alternative (ii) of Lemma~\ref{asc:relative-model}. Over each
vertex $b$ of the ascending Bass--Serre tree $\mathcal B$ place a copy
$T_b$ of $T$. Orient $\mathcal B$ toward its distinguished end, so every
vertex has one parent and any two upward rays eventually agree. The
attaching map to the parent is the homothety $f$. Rescale the metric in
each copy according to its height so that these maps are isometric
embeddings. Their convex directed union is a real tree $\mathcal T$;
it need not be complete. The torus acts on it by homotheties, whose
scales are powers of $\lambda$ determined by height.

\begin{lemma}\label{asc:point-stabilizers}
The height-zero subgroup of $J_\phi$ has trivial stabilizers of
nondegenerate arcs of $\mathcal T$. There are finitely many
$J_\phi$-orbits of points with nontrivial height-zero stabilizer.
For each such point $x$, the full stabilizer $A_x$ has a finite-index
subgroup that is an ascending torus on a proper free factor of $F$.
If $\lambda>1$, the family of these point stabilizers is malnormal.
\end{lemma}

\begin{proof}
For a height-zero element and a compact arc, choose a sufficiently high
ancestor copy containing both endpoints and whose base group contains
the element. A nontrivial element fixing that arc would fix an edge of
the free splitting, which is impossible.

Suppose a nontrivial base element fixes $x$. Its own copy $T_b$ is a
complete simplicial tree, since its finitely many edge lengths have a
positive minimum. It is thus a closed convex subtree of $\mathcal T$.
If $x$ were outside it, the element would fix the segment from $x$ to
its nearest point in $T_b$, contradicting the preceding paragraph.
Hence $x\in T_b$ and is a vertex with nontrivial vertex stabilizer.

Let $C_{b,x}$ be the stabilizer of $x$ in the base group at $b$, and
consider the vertices $b$ for which $C_{b,x}\ne1$. They form a connected
parent-closed subtree of $\mathcal B$: nontrivial elements remain
nontrivial on passing to a parent, and upward rays coalesce. There are
only finitely many types of pairs $(b,x)$ modulo $J_\phi$, because the
free splitting has finite quotient. This proves finiteness of point
orbits. On an upward ray two pair types repeat, yielding an element
$s\in A_x$ of positive height that takes one vertex $b$ to an ancestor.
Its positive translates are cofinal along that ray. Every height-zero
element of $A_x$ belongs to a sufficiently high base group and fixes
$x$ there. Consequently $\ker(h|_{A_x})$ is the increasing union of the
corresponding translates of $C_{b,x}$. The subgroup with heights in
$h(s)\Z$ is therefore the ascending torus based on $C_{b,x}$, with stable
letter $s$ and the appropriate orientation. It has finite index in
$A_x$. The group $C_{b,x}$ is a proper free factor, so has smaller rank.

Finally suppose $\lambda>1$. An element fixing two distinct points has
scale one, hence height zero. It must then be trivial by the arc
assertion. Thus distinct point stabilizers intersect trivially, including
conjugate stabilizers of distinct points. This is precisely malnormality
for a list of orbit representatives.
\end{proof}

If $\lambda=1$, irreducibility and Perron eigenvalue one imply that the
transition matrix is a cyclic permutation matrix. The maps take edges
to single edges. The union tree is simplicial and its torus quotient is
finite. Its edge stabilizers inject into the height group and are
therefore cyclic or trivial. Vertex stabilizers either inject into
height or are those in Lemma~\ref{asc:point-stabilizers}. Cyclic edge
groups are quasiconvex in the hyperbolic torus. The vertex groups are
finitely generated and quasiconvex as well: replace excursions of an
ambient path outside a vertex space by paths in its incident edge
cosets, using their ambient undistortion. Induction applies to the
smaller-rank ascending tori, and the quasiconvex hierarchy theorem
\cite{Wise2021} proves Theorem~\ref{asc:vcs} in this case.

We henceforth assume $\lambda>1$.

\begin{lemma}\label{asc:quasiconvex-peripherals}
Each subgroup $A_x$ in Lemma~\ref{asc:point-stabilizers} is quasiconvex in
$J_\phi$. Consequently the finite malnormal family of these subgroups
is a relatively hyperbolic peripheral structure. If the inductive
hypothesis holds below rank $\rk F$, every member is hyperbolic and
virtually compact special.
\end{lemma}

\begin{proof}
Realize the free splitting as a finite graph with a bouquet at each
vertex carrying its vertex free factor; its other edges realize the
edges of the splitting. Represent $f$ by a map of this graph that sends
bouquets into bouquets and the other edges along their prescribed tree
paths, with bouquet words between successive edges. Such a representative
can be constructed equivariantly on universal covers: choose the
prescribed stabilizer map on each vertex fiber and join the finitely
many representative edge endpoints along their tree paths. The finitely
many bouquet legs in these edge images have a uniform length bound.

Use this representative to form the ordinary tree of graph spaces for
the torus. Its one-skeleton $X$ is a proper cocompact model for the torus
word metric. In every row with $C_{b,x}\ne1$, let $W_b$ be the convex
bouquet lift representing $x$. Let $Y_x$ be the union of these lifts and
the upward strips joining them. It is connected. Finiteness of the pair
types in the proof of Lemma~\ref{asc:point-stabilizers} shows that $A_x$
acts cocompactly on $Y_x$.

Project each retained horizontal row to its convex subtree $W_b$.
These projections are $1$-Lipschitz within rows. They agree, up to a
uniform error, across an upward step. Indeed take the last relative
edge on the path toward the source bouquet. Its image is noncollapsed;
alignment preservation prevents relative backtracking. The target gate
can differ from the image of the source gate only by a terminal bouquet
leg in the image of that one edge. This leg has the uniform bound just
specified. Thus the projections are uniformly Lipschitz across all
vertical edges between retained rows, for the intrinsic metric of $Y_x$.

It remains to treat a component of omitted rows. It attaches to the
retained subtree of $\mathcal B$ by one step $d\to b$ with $C_{d,x}=1$.
The gates to $W_b$ of the entire immediate attaching image have uniformly
bounded diameter. To verify this, use the embedding of relative trees.
If that image misses $x$, its connected image has one gate. If it meets
$x$ at an interior point of a relative edge, only the image of that one
edge can meet the bouquet, and its bouquet legs have bounded length.
If it meets $x$ at a vertex, that vertex has trivial stabilizer: otherwise
its injected stabilizer would give $C_{d,x}\ne1$. Such a vertex has
uniformly finite valence, and only its point image and the finitely many
bounded initial edge legs can meet $W_b$. In each case the actual
attaching image is connected in the horizontal tree. Every gate from
that image to a bouquet that it meets lies in their intersection, so
the preceding bounds apply to the whole image, not merely its vertices
near the attachment.

Map the entire omitted component to a point within bounded distance of
these gates. The resulting map on the vertices of $X$ is a coarse
Lipschitz retraction to $Y_x$: horizontal edges, retained vertical edges,
and the boundary of every omitted component have just been checked.
Therefore inclusion of $Y_x$ with its intrinsic path metric is a
quasi-isometric embedding. Cocompactness of the $A_x$-action proves that
$A_x$ is undistorted, hence quasiconvex in the hyperbolic group $J_\phi$.

A finite malnormal family of quasiconvex subgroups of a hyperbolic group
is a relatively hyperbolic peripheral structure \cite{Bowditch2012}.
Each $A_x$ is itself hyperbolic. Its finite-index ascending torus has
smaller-rank base by Lemma~\ref{asc:point-stabilizers}; induction and
finite-index invariance give the final assertion.
\end{proof}

\subsection{Deep fillings with train-track finite covers}

\begin{lemma}\label{asc:good-train-track-filling}
Assume the inductive hypothesis below rank $\rk F$ and use the peripheral
family in Lemma~\ref{asc:quasiconvex-peripherals}. For every prescribed
finite avoidance set there is a peripheral filling
\[
 \bar J=J_\phi/\langle\!\langle N_x\rangle\!\rangle
\]
that is hyperbolic, has finite-by-cyclic peripheral images, and has a
finite-index subgroup equal to the mapping-torus group of an expanding
train-track map $p:V\to V$ of a finite connected graph. After passing to
a height power, the edge-transition components of $p$ are disjoint
primitive irreducible blocks with connected invariant supports.
\end{lemma}

\begin{proof}
Choose a finite-index normal subgroup $H\lhd J_\phi$ and put
\[
 N_x=H\cap\ker h\cap A_x.
\]
Residual finiteness permits $H$ to avoid any prescribed finite set of
nontrivial elements, simultaneously for all the finitely many
peripherals. The image of $\ker(h|_{A_x})$ in $A_x/N_x$ is finite, whereas
its height image is cyclic. Thus $A_x/N_x$ is finite-by-cyclic.
For sufficiently deep avoidance, the relative filling theorem
\cite{Osin2007} gives exact peripheral reduction and a relatively
hyperbolic quotient with these peripheral images. As the peripheral
images are hyperbolic, $\bar J$ is hyperbolic.

Let $C_v$ be the vertex groups of the base free splitting. Define
$\bar F$ by replacing them by $C_v/(C_v\cap H)$, keeping the underlying
graph and trivial edge groups. Equivalently, quotient $F$ by the normal
closure of all its vertex kernels. The induced tree is denoted
$T^{\mathrm{fin}}$; it is the Bass--Serre tree of this graph of finite
groups. Every $N_x$ is exhausted by the vertex kernels along its upward
union description in Lemma~\ref{asc:point-stabilizers}. Consequently
$\bar J$ has the presentation torus of the induced endomorphism of
$\bar F$. This is a statement about presentations and does not assert
that this endomorphism is injective. The relative tree map descends to
an equivariant edge-path map of $T^{\mathrm{fin}}$.

We impose one additional finite avoidance condition that makes every
iterate of every edge reduced in $T^{\mathrm{fin}}$. There are finitely
many seed turns, namely the internal turns of the images of representative
edges. All turns in later iterates are translates of derivatives of
these seed turns. Here the derivative sends an edge germ to the initial
germ of its image. Before filling their two directions are distinct,
because $f$ is an embedding.

For each seed turn whose directions ever enter the same oriented edge
orbit, choose the first such derivative. Its two directions differ by
a nonidentity multiplier in their vertex stabilizer. Add that multiplier
to the finite avoidance set for $H$. At subsequent derivatives the two
directions stay in the same oriented edge orbit, and their multipliers
are obtained from the first one by $\phi$ and changes of vertex frame.
Since $\phi$ is conjugation in $J_\phi$ and $H$ is normal there, none of
these multipliers belongs to $H$. Directions in different oriented edge
orbits cannot identify under vertex reduction; directions in the same
orbit identify precisely when their multiplier is in the vertex kernel.
Thus no derivative of a seed turn folds. This proves the assertion for
all edge iterates at once.

Put $D=F/(F\cap H)$. Conjugation by the stable letter induces an
automorphism $\sigma$ of this finite group: its induced injection into
itself is necessarily surjective. The natural map $\bar F\to D$ is
injective on every finite vertex group. Its kernel acts freely on
$T^{\mathrm{fin}}$ and has a finite connected graph quotient $V$.
The descended map gives a train-track map $p:V\to V$, satisfying
\[
 p(dz)=\sigma(d)p(z)\qquad(d\in D).
\]
All edges expand, since their relative lengths are multiplied by
$\lambda>1$. With compatible basepoint and stable-letter markings, its
mapping torus is a finite cover of the presentation torus for $\bar F$:
the map on the finite fiber is the permutation $\sigma$.
Equivalently it is the preimage of the stable cyclic subgroup under
$\bar J\to D\rtimes_\sigma\Z$, and has index $|D|$.
This argument remains valid without injectivity of the endomorphism of
$\bar F$ or of $p_*$.

Finally consider reachability of unoriented edges in $V$. If $\alpha$
reaches $\beta$, irreducibility downstairs gives a path of transitions
from $\beta$ to $d\alpha$ for some $d\in D$. The composite gives an
occurrence $\alpha\to d\alpha$ in an iterate of length $m$.
Repeated substitution follows the permutation
\[
 D\longrightarrow D,\qquad s\longmapsto\sigma^m(s)d.
\]
Its orbit of the identity closes. The resulting closed chain of
occurrences passes through $\beta$, so $\beta$ reaches $\alpha$.
Every transition is therefore contained in a strongly connected
component, with no arrows between distinct components. Take a power to
split their periods, making all blocks primitive. A sufficiently long
iterate of any one edge contains all edges of its block and contains
no edges outside it; hence each block support is connected. It is
invariant and has a periodic vertex. Taking a further power fixes all
periodic vertices. This gives the claimed form.
\end{proof}

\subsection{Cubulating the filled groups}

The next lemma treats the possible failure of fundamental-group
injectivity in the preceding finite graph model.

\begin{lemma}\label{asc:injective-block}
Let $p:V_0\to V_0$ be a based train-track map of a finite connected graph,
fixing its base vertex. Suppose all edges expand, the transition matrix
is primitive, and its mapping-torus group is hyperbolic. Then that group
is virtually compact special, whether or not $p_*$ is injective.
\end{lemma}

\begin{proof}
The free groups $P_m=p_*^m(\pi_1V_0)$ have nonincreasing ranks. Choose
$m\ge1$ after these ranks stabilize. The surjection
$p_*:P_m\to P_{m+1}$ is between free groups of the same finite rank, so
is injective by Hopficity. Thus its restriction to $P_m$ is an injective
endomorphism.

Let $\rho:\widehat V\to V_0$ be the based cover corresponding to $P_m$.
Lift $p^m$ to $q:V_0\to\widehat V$ and lift $p$ to
$\widehat p:\widehat V\to\widehat V$, with both maps fixing the chosen
basepoint as appropriate. Uniqueness of based lifts gives
$\widehat p q=qp$. The image graph $W=q(V_0)$ is finite and connected,
and carries exactly $P_m$: its inclusion into the cover is injective
on fundamental groups, while the image of $q_*$ already projects onto
all of $P_m$. The restriction $p_W=\widehat p|_W$ is thus injective on
fundamental groups. It is a train-track map and every edge expands,
because projection takes its edge iterates to the corresponding reduced
expanding iterates of $p$.

The maps $q:V_0\to W$ and $\rho|_W:W\to V_0$ intertwine $p$ and $p_W$.
Their composites are $p^m$ and $p_W^m$: the latter equality again follows
from uniqueness of based lifts. Lemma~\ref{asc:shift-equivalence}
identifies their mapping-torus groups.

We check irreducibility of $p_W$, since mere restriction to a finite
image graph does not by itself establish it. Given edges $e',f'$ of
$W$, choose edges $e,f$ of $V_0$ whose $q$-images contain them.
Primitivity supplies a fixed $k$ such that the image under $p^k$ of each
edge contains $f$ or its reverse. Consequently every path
$q p^k(a)$, for an edge $a$ of $V_0$, contains $q(f)$ or its reverse.
Their lengths have a common finite bound $B$. For $n\ge k$ the reduced
path $q p^n(e)$ is a concatenation of these bounded chunks. Reduction
causes no cancellation, since its projection is the train-track path
$p^{n+m}(e)$. The path $p_W^n(e')$ is a subpath of $q p^n(e)$ and its
length tends to infinity. Once its length exceeds $2B$, it contains
an entire chunk and hence an occurrence of $f'$. Thus every edge
reaches every other edge.

The map $p_W$ satisfies exactly the hypotheses of
\cite[Theorem~B]{HagenWise2015}: it is a fundamental-group-injective
train-track map of a finite graph, all edges expand, its transition
matrix is irreducible, and its mapping-torus group is hyperbolic.
That theorem gives a free cocompact cubulation. Agol's theorem
\cite{Agol2013} gives virtual compact specialness.
\end{proof}

\begin{lemma}\label{asc:filled-vcs}
The fillings in Lemma~\ref{asc:good-train-track-filling} are virtually
compact special.
\end{lemma}

\begin{proof}
Use the finite-index graph torus supplied there, and pass to the height
power for which all its blocks are primitive and its periodic vertices
are fixed. At every vertex of $V$ make a separate copy in each incident
block, retain one central vertex, and join the central vertex to its
block copies by connector intervals. Collapsing each connector star
recovers $V$, so this operation is a homotopy equivalence. The map acts
on each block as before, sends central vertices according to $p$, and
sends each connector linearly to its corresponding connector.

The eventual image of this graph map is connected and contains all
blocks, since each block map is primitive. Its remaining intervals and
central vertices are periodic. Restriction to the eventual image does
not change the mapping-torus group: inclusion and a sufficiently high
power of the map satisfy Lemma~\ref{asc:shift-equivalence}. Connector
orbits now give annuli joining the point tori and the block tori.
This is a finite graph-of-groups decomposition over infinite cyclic
groups. Each edge injection is detected by height, so is injective.

The total group is hyperbolic, as a finite-index subgroup of $\bar J$.
Its cyclic edge groups are quasiconvex, and the finitely generated
vertex groups are consequently quasiconvex. In particular every block
torus is hyperbolic. Each block has a fixed vertex and satisfies the
other hypotheses of Lemma~\ref{asc:injective-block}, so is virtually
compact special. The point tori are infinite cyclic. The quasiconvex
hierarchy theorem \cite{Wise2021} now gives virtual compact specialness
of the total group, and therefore of its finite-index supergroup
$\bar J$.
\end{proof}

\begin{proof}[Proof of Theorem~\ref{asc:vcs}]
The base cases and the automorphism case were disposed of at the start.
Apply Lemma~\ref{asc:relative-model}. Its rank-reduction alternative is
settled by induction and finite-index invariance. In the tree-model
alternative, the case $\lambda=1$ was settled following
Lemma~\ref{asc:point-stabilizers}. For $\lambda>1$,
Lemma~\ref{asc:quasiconvex-peripherals} supplies a finite malnormal
quasiconvex family of hyperbolic virtually compact special subgroups.
Lemmas~\ref{asc:good-train-track-filling} and \ref{asc:filled-vcs}
produce arbitrarily deep fillings to hyperbolic virtually compact
special groups, with their inherited finite-by-cyclic peripheral images.
Finite-by-cyclic groups are virtually cyclic and hence virtually compact
special. Finally Lemma~\ref{asc:field-hypotheses} verifies the coherence,
local-indicability, finiteness, dimension and field-acyclicity assumptions
of Theorem~\ref{fill:criterion}. That criterion applies and proves that
$J_\phi$ is virtually compact special. The induction is complete.
\end{proof}

\begin{corollary}\label{cor:ascending-finite-rank}
Let $F$ be a free group of finite rank and let $\phi:F\to F$ be injective.
If $J_\phi$ is word-hyperbolic, it has a finite-index subgroup $H$ fitting
into an exact sequence
\[
 1\longrightarrow F_d\longrightarrow H\longrightarrow\Z
 \longrightarrow1
\]
for some finite $d\geq0$, where $F_d$ is the rank-$d$ free group.
Thus $H\cong F_d\rtimes_\alpha\Z$ for an automorphism $\alpha$ of $F_d$.
\end{corollary}

\begin{proof}
Theorem~\ref{asc:vcs} makes $J_\phi$ virtually RFRS.
Lemma~\ref{asc:field-hypotheses} supplies a finite classifying complex
whose chains are acyclic over the subgroup division field.
Corollary~\ref{fill:rfrs-novikov} therefore gives two-sided ordinary
Novikov contractions on a finite-index subgroup $H$. The restricted
group-ring complex is still the cellular complex of a connected
universal cover, so its ordinary $H_0$ is $\Q$. Its character is
nonzero: a zero character would make the Novikov ring $\Q H$ and the
contraction would contradict this $H_0$.

Lemma~\ref{fill:bns} gives an epimorphism $H\to\Z$ with finitely
generated kernel. Coherence and cohomological dimension at most two
pass to $H$ by Lemma~\ref{asc:field-hypotheses}. The kernel is therefore
finitely presented, and Lemma~\ref{fill:free-kernel} makes it free.
Its rank is finite by finite generation. A lift of $1\in\Z$ splits
the sequence. This proof applies after every branch of the induction
establishing Theorem~\ref{asc:vcs}.
\end{proof}

The epimorphism in this corollary need not be the restriction of the
canonical height homomorphism of the original ascending presentation.
